\documentclass[../../../include/open-logic-section]{subfiles}

\begin{document}

\olfileid[hi]{sth}{cardinals}{cp}
\olsection{कैंटर का सिद्धांत}

\olref[ordinals][vn]{sec} को याद कीजिए। हम सुव्यवस्थित समुच्चयों पर चर्चा
कर रहे थे और हमने सुझाया था कि सुव्यवस्थाओं के प्रतिनिधि के रूप में काम
करने वाली वस्तुएँ होना उपयोगी होगा। इसी विचार से हमने क्रमसूचक प्रस्तुत किए
और फिर \olref[ordinals][ordtype]{ordtypesworklikeyouwant} में दिखाया कि वे
हमारी अपेक्षा के अनुसार व्यवहार करते हैं, अर्थात्:
\[
	\ordtype{A, <} = \ordtype{B, \lessdot} 
	\text{ तभी जब } \tuple{A, <} \isomorphic \tuple{B, \lessdot}.
\]
और पीछे \olref[sfr][siz][equ]{sec} को याद कीजिए। वहाँ सहज ढंग से काम करते
हुए हमने किसी समुच्चय के ``आकार'' की धारणा प्रस्तुत की थी। विशेषतः हमने कहा
था कि दो समुच्चय समसंख्य हैं, $\cardeq{A}{B}$, ठीक तभी जब कोई
!!a{bijection} $f \colon A \to B$ हो। यह धारणा सुव्यवस्था की धारणा से
स्वभावतः {सरल} है: हमें केवल !!{bijection}ों में रुचि है, इसमें नहीं कि वे
!!{bijection} ``किसी संरचना को सुरक्षित रखते हैं'' या नहीं (जैसा
क्रम-समाकृतिकताओं में होता है)।

इन सब से एक सहज विचार उत्पन्न होता है। जैसे हमने सुव्यवस्थाओं को मापने के
लिए कुछ वस्तुएँ, \emph{क्रमसूचक}, प्रस्तुत की थीं, वैसे ही आकार मापने के
लिए कुछ वस्तुएँ, \emph{कार्डिनल}, प्रस्तुत कर सकते हैं। यही इस अध्याय का
उद्देश्य है।

ये कार्डिनल क्या होंगे यह कहने से पहले हमें वह सिद्धांत निर्धारित करना
चाहिए जिसका उन्हें पालन करना होगा। समुच्चय $X$ की कार्डिनलता को $\card{X}$
लिखते हुए हम चाहेंगे कि वे इसका पालन करें:
\[
	\card{A} = \card{B} \text{ तभी जब } \cardeq{A}{B}.
\]
हम इसे \emph{कैंटर का} सिद्धांत कहेंगे, क्योंकि शायद कैंटर ही पहले व्यक्ति
थे जिनके मन में यह बिल्कुल स्पष्ट रूप से था। (\olref[hp]{sec} में हम
\emph{ह्यूम के} सिद्धांत से इसके संबंध पर और चर्चा करेंगे।) अतः हमारा
उद्देश्य प्रत्येक $X$ के लिए $\card{X}$ को इस प्रकार परिभाषित करना है कि
उससे कैंटर का सिद्धांत प्राप्त हो।

\end{document}
